\documentclass[11pt]{book}
\usepackage{cours-preambule}

\toggletrue{student}
\HideSolutionstrue

\makeatletter
\renewcommand{\@chapapp}{MPSI3 -- \today\ -- Devoir surveillé}
\makeatother

\renewcommand{\thesection}{E\arabic{section}}
\titleformat{\section}{\Large\bfseries}{
    \begin{tikzpicture}[scale=0.55, overlay]
        \draw[line width=1pt, color=darkgray, xshift=1em, yshift=0.6em]
            (1,-0.4) -- (1,0.8) --(0.8,1) -- (-0.8,1) -- (-1,0.8) -- (-1,-0.8) --
            (-0.8,-1) -- (33,-1);
        \node[xshift=0.6em, yshift=0.35em] at (0,0)
            {\rmfamily\textcolor{darkgray}{\thesection}};
    \end{tikzpicture}}{0.4em}{\hspace{1.4em}}

\setlist[enumerate,1]{leftmargin=10pt, label=\sqenumi}
\setlist[enumerate,2]{leftmargin=20pt, label=\Alph*)}

\newcommand{\figsvg}[1]{
  \begin{center}
    \subimport{figures/}{#1}
  \end{center}
}

\begin{document}
\setcounter{chapter}{9}
\chapter{Cristallographie et induction \ifstudent{-- corrigé}}
\label{ch:ds10}

\ifprof{
\begin{center}
	\Large\bfseries
	\ul{Tout moyen de communication est interdit}
	\smallbreak
	\ul{Les téléphones portables doivent être éteints et rangés dans les sacs}
	\smallbreak
	\ul{Les calculatrices sont \textit{interdites}}
\end{center}
Le devoir est composé des exercices \textit{indépendants} suivants~:
{\Large
\begin{itemize}[label=$\diamond$]
	\litem{Exercice 1}~: Oxyde de zirconium solide
	\litem{Exercice 2}~: Questions de cours
	\litem{Exercice 3}~: Rails de \textsc{Laplace} inclinés
\end{itemize}
}

Les différentes questions peuvent être traitées dans l'ordre désiré.
\textbf{Cependant}, le numéro complet de la question doit être indiqué, et
\textbf{vous indiquerez si vous traitez la question d'un exercice sur une page
	complètement déconnectée}, sous peine de n'être ni vue ni corrigée.
\bigbreak
Une attention particulière sera portée à la \textbf{qualité de rédaction}. Les
hypothèses doivent être clairement énoncées, les propositions reliées entre
elles par des connecteurs logiques, les lois et théorèmes énoncés, sans pour
autant devenir une composition de français.
\bigbreak
De plus, la \textbf{présentation} de la copie sera prise en compte. Outre la
numérotation des questions, l'écriture, l'orthographe, les encadrements, la
marge, le cadre laissé pour la note et le commentaire font partie des points à
travailler. Il est notamment attendu que \textbf{les expressions littérales
	soient encadrées}, que \textbf{les calculs n'apparaissent pas} mais que le
détail des grandeurs avec leurs unités soit indiqué, et \textbf{les applications
	numériques soulignées}.
\bigbreak
Ainsi, l'élève s'expose aux malus suivants concernant la forme et le fond~:
\begin{tprop}{Malus, dngr}
	\begin{minipage}{0.50\linewidth}
		\begin{itemize}
			\item A~: application numérique mal faite~;
			\item N~: numéro de copie manquant~;
			\item P~: prénom manquant~;
			\item E~: manque d'encadrement des réponses~;
			\item M~: marge non laissée ou trop grande~;
			\item V~: confusion ou oubli de vecteurs~;
		\end{itemize}
	\end{minipage}
	\begin{minipage}{0.50\linewidth}
		\begin{itemize}
			\item Q~: question mal ou non indiquée (même passées)~;
			\item C~: copie grand carreaux~;
			\item U~: mauvaise unité (flagrante)~;
			\item H~: homogénéité non respectée~;
			\item S~: chiffres significatifs non cohérents~;
			\item $\f$~: loi physique fondamentale brisée.
		\end{itemize}
	\end{minipage}
\end{tprop}

\begin{timpo}{Exemple application numérique, hand}
	\vspace*{-10pt}
	\begin{minipage}{0.45\linewidth}
		\begin{gather*}
			\boxed{n = \frac{PV}{RT}}
			\qav
			\left\{
			\begin{array}{rcl}
				p & = & \SI{1.0e5}{Pa}                \\
				V & = & \SI{1.0e-3}{m^3}              \\
				R & = & \SI{8.314}{J.mol^{-1}.K^{-1}} \\
				T & = & \SI{300}{K}
			\end{array}
			\right.\\
			\mathrm{A.N.~:}\quad
			\ul{n = \SI{5.6e-4}{mol}}
		\end{gather*}
	\end{minipage}
	\hfill
	\cancel{\bcancel{
			\begin{minipage}{0.45\linewidth}
				\begin{gather*}
					n = \frac{PV}{RT} = \frac{\num{e5}\cdot\num{1}}{8.32\cdot300}
					= 0.56
				\end{gather*}
			\end{minipage}
		}}
\end{timpo}
\newpage
}

\section{Oxyde de zirconium solide \hfill \textcolor{red}{/30}}
\label{sec:oxzir}

\switch{
	\begin{center}
		\begin{framed}
			\Large\bfseries
			Aucune impasse sur la cristallographie ne sera tolérée.
		\end{framed}
	\end{center}

	Les piles à combustible à oxyde solide permettent d'avoir en contact deux phases
	solide et gazeuse, ce qui supprime les problèmes liés à la gestion de trois
	phases, notamment la corrosion. Les électrodes sont poreuses de façon à
	permettre un transport rapide des gaz. Un matériau de choix pour l'électrolyte
	est l'oxyde de zirconium, appelé zircone, stabilisé à l'yttrium.
	\smallbreak
	La zircone peut être assimilée à un cristal ionique formé de cations
	\ce{Zr^{4+}} et d'anions \ce{O^{2-}} de rayons respectifs $r_+$ et $r_-$. Les
	cations sont distribués aux nœuds d'un réseau cubique faces centrées (CFC).
	\bigbreak
	\begin{enumerate}
		\item Rappeler ce qu'est le modèle des sphères dures.
		\item Représenter la maille conventionnelle d'une structure CFC de cations.
		      Indiquer le nombre de cations en propre par maille en détaillant le
		      calcul.
		\item Définir puis donner sans démonstration la compacité d'une telle
		      structure dans le cas d'une maille métallique. Commenter.
		\item Indiquer et représenter où se situent les sites tétraédriques de cette
		      maille. Combien y en a-t-il~?
		\item Exprimer le rayon maximal $r_-$ de la particule sphérique pouvant
		      s'insérer dans ces sites sans induire de déformation en fonction du
		      paramètre de maille $a$ et de $r_+$. La condition de contact sera
		      énoncée en français et représentée schématiquement.
	\end{enumerate}
	\textbf{Les anions occupent tous les sites tétraédriques de la maille CFC formée
		par les cations.}
	\begin{enumerate}[resume]
		\item Déterminer le nombre d'anions contenus dans cette maille.
		\item En déduire la formule de la zircone.
		\item Définir et donner, en justifiant, la coordinence des anions par rapport
		      aux cations, et des cations par rapport aux anions.
		\item Exprimer la masse volumique de la zircone en fonction de $a$, de
		      $M_{\ce{Zr}}$ la masse molaire du zirconium, de $M_{\ce{O}}$ la masse
		      molaire de l'oxygène et du nombre d'\textsc{Avogadro}.
	\end{enumerate}
	\textbf{La formule de l'oxyde d'yttrium est \ce{Y2O3}.}
	\begin{enumerate}[resume]
		\item En déduire la charge du cation yttrium.
		\item Le dopage consiste à substituer dans la maille élémentaire de l'oxyde de
		      zirconium une fraction molaire $x$ des cations \ce{Zr^{4+}} par des
		      cations yttrium. Expliquer pourquoi l'électroneutralité de la structure
		      n'est alors pas respectée.
		\item Proposer une modification de la formule chimique impliquant le nombre
		      $y$ d'anions \ce{O^{2-}} présents dans la zircone dopée à l'oxyde
		      d'yttrium, au moyen de $x$, pour rétablir cette électroneutralité.
	\end{enumerate}
	\newpage
}{
	\begin{enumerate}
		\nitem{1} Le modèle des sphères dures consiste à considérer les entités
		monoatomiques (atomes ou ions) comme des sphères impénétrables et
		indéformables.
	\end{enumerate}
	\vspace*{-10pt}
	\noindent
	\begin{minipage}[t]{.73\linewidth}
		\begin{enumerate}[start=2]
			\nitem{3} On a 8 atomes aux sommets, et 6 atomes au centre des faces (voir
			Figure~\ref{fig:cfc}). Les premiers sont partagés entre huit mailles, les
			seconds entre 2 mailles, d'où
			\[
				N = 6 \times \frac{1}{2} + 8 \times \frac{1}{8}
				\Lra
				\boxed{N = 4} \text{ cations par maille}
			\]
			\nitem{3} La compacité est le rapport du volume occupé par les entités sur
			le volume d'une maille. Pour une CFC, elle est de \textbf{74\%}. C'est la
			compacité maximale pour une structure de sphères dures identiques. Il reste
			donc \textbf{26\% de vide} qui peut être occupé par d'autres entités plus
			petites.
			\nitem{3} Dans la CFC, les sites tétraédriques (sites T) sont \textbf{au
				centre des huit cubes d'arête $a/2$}. Il y en a donc \textbf{8 par maille}
			(voir Figure~\ref{fig:sitesT}).
			\nitem{4} À la limite où l'anion occupe le site T sans déformer la structure
			de cations, le contact anion-cation se fait sur la grande diagonale du petit
			cube d'arête $a/2$ (voir Figure~\ref{fig:habsiteT}), soit
			\[
				r_+ + r_- \leq \frac{1}{2}\times \frac{a \sqrt{3}}{2}
				\Lra
				\boxed{r_- \leq \frac{a \sqrt{3}}{4} - r_+}
			\]
			\nitem{1} Les sites T appartiennent en propre à la maille, il y a donc autant
			d'anions que de site, c'est-à-dire \textbf{8 anions par maille}.
			\nitem{2} Il y a donc deux fois plus d'anions \ce{O^{2-}} que de cations
			\ce{Zr^{4+}}. La formule avec les plus petits entiers naturels est donc
			\ce{ZrO2}.
			\nitem{3} La coordinence d'une entité est son nombre de plus proches voisins
			(PPV). Ici, un anion est dans un site T, donc entouré de 4 cations PPV, donc
			la coordinence des \textbf{anions par rapport aux cations est de 4}. Un cation
			au centre d'une face, lui, est entouré de 8 petits cubes d'arête $a/2$, au
			centre desquels se trouvent les anions~: la coordinence des \textbf{cations
				par rapport aux anions est de 8}.
		\end{enumerate}
	\end{minipage}
	\hfill
	\begin{minipage}[t]{.25\linewidth}
		~
		\begin{center}
			\includegraphics[scale=.8]{cfc}
			\captionof{figure}{Maille d'une CFC.}
			\label{fig:cfc}
		\end{center}
		\begin{center}
			\includegraphics[scale=1]{sitesT}
			\captionof{figure}{Sites T de la structure CFC.}
			\label{fig:sitesT}
		\end{center}
		\begin{center}
			\includegraphics[scale=1]{habsiteT}
			\captionof{figure}{Habitabilité d'un site T d'une CFC.}
			\label{fig:habsiteT}
		\end{center}
	\end{minipage}
	\begin{enumerate}[start=9]
		\nitem{3} La masse volumique est le rapport de la masse des entités \textit{en
			propre} sur le volume de la maille. Ainsi,
		\[
			\rho = \frac{N_{\ce{Zr}} m_{\ce{Zr}} + N_{\ce{O}} m_{\ce{O}}}{a^3}
			\Lra
			\boxed{\rho = \frac{4M_{\ce{Zr}} + 8M_{\ce{O}}}{\Nc_{\rm A}a^3}}
		\]
		\nitem{2} L'oxygène se trouve toujours sous la forme d'anions \ce{O^{2-}}. Par
		neutralité de la structure, on a donc
		\[
			2z + 3 \times (-2) = 0
			\Lra
			z=3
		\]
		d'où la notation \ce{Y^{3+}}.
		\nitem{2} Si la substitution se fait le plus simplement possible, un cation
		\ce{Zr^{4+}} est remplacé par un cation \ce{Y^{3+}}~: tout se passe donc comme
		si le cristal avait perdu une charge $+$, et donc qu'il s'est chargé
		négativement.
		\nitem{3} En fonction de $x$ et $y$, la formule chimique de la structure
		s'écrit $\ce{Zr_{1-x}Y_{x}O_{y}}$. La neutralité impose
		\[
			4(1-x) + 3x -2y = 0
			\Lra
			\boxed{y = 2 - \frac{x}{2}}
		\]
		Ainsi, au cours du processus de substitution, certains sites tétraédriques de
		l'alliage se vident de leurs ions \ce{O^{2-}} pour préserver la neutralité de
		la structure. On vérifie par ailleurs qu'avec $x=0$, c'est-à-dire sans dopage,
		on retrouve bien $y=2$.
	\end{enumerate}
}
\vspace*{-10pt}
\section{Questions de cours \hfill \textcolor{red}{/36}}
\label{sec:cours}
\switch{
	\begin{center}
		\begin{framed}
			\large\bfseries
			Chaque sous-exercice est indépendant des autres.
		\end{framed}
	\end{center}

	\subsection{Couple sur un aimant}
	\label{ssec:cplaimant}
	Soit un aimant de moment magnétique $\vv{m}$ plongé dans un champ magnétique
	$\vv{B}$ dans le plan de ce moment magnétique. On appelle $\tt$ l'angle orienté
	de $\vv{B}$ à $\vv{m}$.
	\begin{enumerate}
		\item Faire un schéma de la situation, et établir entièrement le système
		      d'étude.
		\item Exprimer le couple de \textsc{Laplace} subi par $\vv{m}$ en fonction
		      de $\tt$.
		\item En déduire les positions d'équilibre de $\vv{m}$.
		\item En étudiant dans quel sens le couple de \textsc{Laplace} tend à faire
		      tourner $\vv{m}$ en cas de petites perturbations, déterminer laquelle des
		      deux positions d'équilibre est stable, et laquelle est instable.
	\end{enumerate}

	\subsection{Modération de \textsc{Lenz}}
	\label{ssec:lenz}
	\noindent
	\begin{minipage}[t]{.47\linewidth}
		Dans chacun des circuits ci-contre, la spire circulaire et/ou l'aimant droit
		sont déplacés dans le sens indiqué par la double flèche.
		\begin{enumerate}[resume]
			\item Énoncer la loi de \textsc{Lenz}.
			\item Indiquer le signe du courant $i$ apparaissant dans la spire pendant le
			      déplacement, en justifiant en français et par un schéma pour chaque
			      situation particulière.
		\end{enumerate}
	\end{minipage}
	\hfill
	\begin{minipage}[t]{.47\linewidth}
		~
		\begin{center}
			\includegraphics[width=\linewidth]{inducsigne}
			\label{fig:inducsigne}
		\end{center}
	\end{minipage}

	\subsection{Loi de \textsc{Faraday}}
	\label{ssec:faraday}
	\noindent
	\begin{minipage}[t]{.75\linewidth}
		On considère un circuit carré de côté $a$ et de résistance totale $R$, situé
		dans un plan orthogonal à un champ magnétique uniforme mais \textbf{variable}
		$\vv{B}(t) = B_0\exr^{-t/\tau}\uz$ avec $B_0$ et $\tau$ strictement positifs.
		\begin{enumerate}[start=7]
			\item Expliquer pourquoi il y a induction.
			\item Exprimer l'intensité $i$ du courant induit représenté sur le schéma.
			\item Vérifier que son signe soit en accord avec la loi de \textsc{Lenz}.
			\item Quel autre paramètre peut-on faire varier pour avoir induction~? Citer
			      un exemple de système présentant de l'induction dans ce cas.
		\end{enumerate}
	\end{minipage}
	\hfill
	\begin{minipage}[t]{.20\linewidth}
		~
		\vspace*{-20pt}
		\begin{center}
			\includegraphics[scale=1]{faraexo}
			\label{fig:faraexo}
		\end{center}
	\end{minipage}

	\subsection{Inductance propre}
	\label{ssec:indprp}
	\begin{enumerate}[start=11]
		\item Définir ce qu'est le flux propre d'un circuit.
		\item Comment se définit l'auto-inductance~?
		\item Calculer l'inductance propre d'une bobine de TP, avec $N$ spires de
		      rayon $a = \SI{3}{cm}$ et de longueur $\ell = \SI{10}{cm}$, avec $\mu_0 =
			      4\pi \cdot \SI{e-7}{H.m^{-1}}$.
	\end{enumerate}

	\newpage
}{
	\small
	\subsection{Couple sur un aimant \hfill \textcolor{red}{/10}}
	\label{ssec:cplaimant}
	\noindent
	\begin{minipage}[t]{.7\linewidth}
		\begin{enumerate}
			\nitem{4} Voir Figure~\ref{fig:aim_cpl}. On étudie l'aimant, caractérisé
			par un moment $\vv{m}$ uniquement, dans un champ magnétique uniforme et
			stationnaire $\vv{B}$ dans le plan de l'aimant. On utilise le référentiel
			du laboratoire supposé galiléen, avec une base cartésienne $(\ux,\uy,\uz)$
			orthonormée directe. On appelle $\tt$ l'angle orienté de $\vv{B}$ à
			$\vv{m}$, compté positivement dans le sens direct.
		\end{enumerate}
	\end{minipage}
	\hfill
	\begin{minipage}[t]{.3\linewidth}
		~
		\vspace*{-15pt}
		\begin{center}
			\includegraphics[scale=1]{aim_cpl}
			\captionof{figure}{Aimant dans $\protect \vv{B}$.}
			\label{fig:aim_cpl}
		\end{center}
	\end{minipage}
	\begin{enumerate}[start=2]
		\nitem{2}\leftcenters{
			\hspace{-12pt}
			On a}{
			$\vv{m} = m \left( \cos{\theta}\ux + \sin{\theta}\uy \right)$
			et
			$\vv{B} = B_0\ux$
		}
		\smallbreak
		\vspace{-40pt}
		\leftcenters{D'où}{$\DS \vv{\Gamma}_z = (\vv{m}\wedge \vv{B})\cdot \uz \Lra
				\boxed{\vv{\Gamma}_z = -mB_0\sin{\theta}\uz}$}
		\nitem{2} On est à l'équilibre si la somme des forces et la somme des moments
		sont nulles. Avec uniquement le couple magnétique, on a bien une
		résultante nulle~: on a donc équilibre si $\Mcf_z = \vv{0}$, soit
		\begin{gather*}
			\sin{\tt_{\rm eq}} = 0
			\qquad \Lra \qquad
			\tt_{\rm eq}       = 0
			\qqou
			\tt_{\rm eq}       = \pi
		\end{gather*}
		\nitem{2}
		\vspace*{-32pt}
		\begin{itemize}[leftmargin=10pt]
			\item En $\tt_{\rm eq} = \pi$, une petite déviation vers le haut
			      donne un mouvement de rotation dans le sens horaire, qui écarte
			      donc l'aimant de sa position d'équilibre~: il est
			      \textbf{instable}.
			\item À l'inverse, en $\tt_{\rm eq} = 0$, une petite
			      déviation vers le haut donne un mouvement de rotation dans
			      le sens direct, le ramenant à sa position d'équilibre~: il
			      est \textbf{stable}.
		\end{itemize}
	\end{enumerate}
	\subsection{Modération de \textsc{Lenz} \hfill \textcolor{red}{/7}}
	\label{ssec:lenz}
	\begin{enumerate}[start=5]
		\nitem{1} \textbf{L'induction modère, par ses conséquences, les causes qui
			lui a donné naissance.}
		\nitem{6} Les lignes de champ vont du Nord vers le Sud. On doit déterminer
		le sens de variation du champ vu par la spire~: le champ induit atténue
		cette variation. On en déduit le signe réel de $i$ par la règle de la main
		droite.
		\smallbreak
		\begin{minipage}[t]{.45\linewidth}
			\begin{enumerate}
				\item $\vv{B}_{\rm aimant}$ augmente vers la gauche~: $\vv{B}_{\rm
						      induit}$ est vers la droite, donc issu de $i_{\rm ind} > 0$.
				      \begin{center}
					      \includegraphics[width=\linewidth]{inducsigne_a}
					      \captionof{figure}{Situation 1.}
					      \label{fig:indsa}
				      \end{center}
				\item Même situation physique, convention opposée~: $i_{\rm ind} < 0$.
				\item Cette fois $\vv{B}_{\rm aimant}$ augmente vers la droite~:
				      $\vv{B}_{\rm ind}$ est vers la gauche. Sens réel du courant opposé~:
				      $i_{\rm ind} < 0$. Voir Figure~\ref{fig:indsb}.
			\end{enumerate}
		\end{minipage}
		\hfill
		\begin{minipage}[t]{.45\linewidth}
			\begin{center}
				\includegraphics[width=\linewidth]{inducsigne_b}
				\captionof{figure}{Situation 3.}
				\label{fig:indsb}
			\end{center}
			\begin{enumerate}[start=4]
				\item Pareil qu'à la question 2~: $i_{\rm ind} < 0$.
				\item Pas de mouvement relatif donc pas de variation de champ donc pas
				      d'induction~: $i_{\rm ind} = 0$.
				\item Mouvement relatif amplifié. Ici, le champ \textbf{diminue} vers la
				      gauche, donc le champ induit le \textbf{renforce}~; ainsi $i_{\rm
							      ind} < 0$.
				      \begin{center}
					      \includegraphics[width=\linewidth]{inducsigne_c}
					      \captionof{figure}{Situation 6.}
					      \label{fig:indsc}
				      \end{center}
			\end{enumerate}
		\end{minipage}
	\end{enumerate}
	\subsection{Loi de \textsc{Faraday} \hfill \textcolor{red}{/11}}
	\label{ssec:faraday}
	\begin{enumerate}[start=7]
		\nitem{4} Le champ magnétique est variable, et la surface $S = a^2$ est fixe.
		Ainsi, le flux est variable. D'après la loi de \textsc{Faraday}, une f.é.m.\
		induite va apparaître dans le circuit~: il y a donc un phénomène
		d'induction.
		\nitem{4}
		On a alors~:
		\vspace{-10pt}
		\begin{gather*}
			\begin{aligned}{}
				\F_S(\vv{B}) & = -Ba^2
				\\
				             & = -B_0a^2\exr^{-t/\tau}
			\end{aligned}
			\qqso
			\begin{aligned}{}
				e_{\rm ind}         & = -\dv{\F}{t}
				\\\Lra
				\Aboxed{e_{\rm ind} & = -\frac{B_0a^2}{\tau}\exr^{-t/\tau} < 0}
			\end{aligned}
		\end{gather*}
		Or, comme le circuit est fermé, \textbf{et que $e$ est en convention
			générateur},
		\vspace{-10pt}
		\begin{align*}
			e_{\rm ind}         & = Ri_{\rm ind}
			\\\Lra
			\Aboxed{i_{\rm ind} & = -\frac{B_0a^2}{T\tau}\exr^{-t/\tau} < 0}
		\end{align*}
		donc l'intensité est \textbf{négative}.
		\nitem{1} En effet, le champ magnétique induit
		réel doit s'opposer à la diminution du champ extérieur $\vv{B}$, en créant un
		champ magnétique positif selon $\uz$~: le sens réel du courant donné par la
		main droite est l'opposé de celui représenté.
		\nitem{2} On peut faire varier la surface. Les rails de \textsc{Laplace}
		générateurs sont un exemple. L'expérience de TP avec le galvanomètre en est
		un autre.
	\end{enumerate}
	\subsection{Inductance propre \hfill \textcolor{red}{/8}}
	\label{ssec:indprp}
	\begin{enumerate}[start=11]
		\nitem{1} Le flux propre, noté $\F_P$ d'un circuit est le flux de son propre
		champ magnétique à travers sa surface.
		\nitem{1} L'auto-inductance est le facteur de proportionnalité entre le flux
		propre et le courant à l'origine de son champ propre~:
		\[
			\F_p(t) = Li(t)
		\]
		\nitem{6} \noindent
		\begin{minipage}[t]{.4\linewidth}
			~
			\begin{center}
				\includegraphics[scale=1]{bobib}
				\captionof{figure}{Champ propre d'un solénoïde.}
				\label{fig:bobib}
			\end{center}
		\end{minipage}
		\hfill
		\begin{minipage}[t]{.69\linewidth}
			Le champ propre dans un solénoïde est
			\[
				\vv{B_p}(t) = \mu_0 \frac{N}{\ell }i(t)\,\uz
			\]

			$i$ et $\vv{B_p}$ respectent la règle de la main droite.
		\end{minipage}
		Pour $N$ spires~:
		\[
			\F_p = N\times \vv{B_p}\cdot \vv{S}
		\]
		Or, $\vv{B_p}$ et $\vv{S}$ sont tous deux orientés à partir de $i$ selon
		la règle de la main droite, donc
		\begin{gather*}
			\vv{S} = S\uz
			\qqet
			\vv{B_p} = \mu_0 \frac{N}{\ell }i(t)\,\uz
			\\\Lra
			\boxed{\F_p = \mu_0 \frac{N^2}{\ell }Si(t)}
		\end{gather*}
		On a démontré que le flux magnétique propre et l'intensité étaient
		proportionnels, et que la constante de proportionnalité était positive.
		On identifie simplement~:
		\[
			\boxed{L = \mu_0 \frac{N^2}{\ell }S}
		\]
		D'où par application numérique
		\[
			\ul{L \approx \SI{35}{mH}}
		\]

	\end{enumerate}
}

\section{Rails de \textsc{Laplace} inclinés \hfill \textcolor{red}{/19}}
\label{sec:rlplincl}

\switch{
	\begin{center}
		\begin{framed}
			\large\bfseries
			Indiquer la ou les bonnes réponses, en justifiant tout votre raisonnement.
			Une réponse sans justification ou avec une justification fausse ne
			rapportera aucun point.
		\end{framed}
	\end{center}

	Soient deux rails rigides, parallèles, distants de $a$, faisant un angle
	$\alpha$ avec le plan horizontal, conducteurs, de résistance négligeable,
	baignant dans un champ magnétique uniforme $\vv{B_0}$ vertical. Les rails sont
	reliés à une source de courant délivrant une intensité $i(t)=I_0\cos(\wt)$. Une
	barre conductrice de résistance négligeable, de masse $m$ peut glisser sans
	frottement sur les rails en restant perpendiculaire aux rails. La barre est
	lâchée sans vitesse initiale depuis la position $x(t=0)=x_0$.

	\figsvg{fig.pdf_tex}

	\begin{enumerate}
		\item Quelle est la résultante des forces de \textsc{Laplace} $\vv{F_L}$
		      s'exerçant sur la barre mobile~?
		      \smallbreak
		      \noindent
		      \begin{minipage}{0.5\linewidth}
			      \begin{enumerate}
				      \item \figsvg{force1.pdf_tex}
				      \item \figsvg{force2.pdf_tex}
			      \end{enumerate}
		      \end{minipage}
		      \begin{minipage}{0.5\linewidth}
			      \begin{enumerate}[start=3]
				      \item \figsvg{force3.pdf_tex}
				      \item \figsvg{force4.pdf_tex}
			      \end{enumerate}
		      \end{minipage}
		\item Donner l'expression de la puissance $\Pc_L$ des actions de
		      \textsc{Laplace}.
		      \smallbreak
		      \noindent
		      \begin{minipage}{0.5\linewidth}
			      \begin{enumerate}
				      \item $P_L=B_0ia \xp$
				      \item $P_L=-B_0ia\xp$
			      \end{enumerate}
		      \end{minipage}
		      \begin{minipage}{0.5\linewidth}
			      \begin{enumerate}[start=3]
				      \item $P_L=-B_0ia\cos(\alpha)\xp$
				      \item $P_L=0$
			      \end{enumerate}
		      \end{minipage}
		\item Quelle est l'équation différentielle vérifiée par $x(t)$~?
		      \smallbreak
		      \noindent
		      \begin{minipage}{0.5\linewidth}
			      \begin{enumerate}
				      \item $m\xpp=mg\sin(\alpha)-B_0ia\cos(\alpha)$
				      \item $m\xpp=mg\cos(\alpha)-B_0ia\cos(\alpha)$
			      \end{enumerate}
		      \end{minipage}
		      \begin{minipage}{0.5\linewidth}
			      \begin{enumerate}[start=3]
				      \item $m\xpp=mg\cos(\alpha)-B_0ia$
				      \item $m\xpp=mg\sin(\alpha)-B_0ia$
			      \end{enumerate}
		      \end{minipage}
		\item Exprimer $x(t)$.
		      \smallbreak
		      \noindent
		      \begin{minipage}{0.5\linewidth}
			      \begin{enumerate}
				      \item $x(t)=\cfrac{1}{2}g\sin(\alpha)t^2+
					            \cfrac{B_0a}{m\w^2}\cos(\alpha) I_0\cos(\wt)+x_0$
				      \item $x(t)=\cfrac{1}{2}g\cos(\alpha)t^2+
					            \cfrac{B_0a}{m}\cos(\alpha) I_0\cos(\wt)+x_0$
			      \end{enumerate}
		      \end{minipage}
		      \begin{minipage}{0.5\linewidth}
			      \begin{enumerate}[start=3]
				      \item $x(t)=\cfrac{1}{2}g\sin(\alpha)t^2+
					            \cfrac{B_0a}{m}\cos(\alpha) I_0(\cos(\wt)-1)+x_0$
				      \item $x(t)=\cfrac{1}{2}g\sin(\alpha)t^2+
					            \cfrac{B_0a}{m\w^2}\cos(\alpha) I_0(\cos(\wt)-1)+x_0$
			      \end{enumerate}
		      \end{minipage}
	\end{enumerate}
}{
	\begin{enumerate}
		\nitem{4} Par définition, ${\rm d}\vv{F_L}=i\vv{\dd{\ell}}\wedge \vv{B_0}$.
		Cette force est orthogonale à la barre et à $\vv{B_0}$. Elle est donc
		horizontale, orientée vers la gauche~: \fbox{\textbf{réponse C)}}
		\figsvg{force3.pdf_tex}
		\nitem{4} Par définition de la puissance d'une force $\Pc_L=\vv{F_L}\cdot\vv{v}$.
		Or $\vv{v}=\xp\ex$.
		\smallbreak
		Donc $\Pc_L=(\vv{F_L}\cdot\ex)\xp$, avec $\vv{F_L}=-iB_0a\vv{n}$ et
		$\vv{n}$ le vecteur unitaire associé à l'axe horizontal vers la droite.
		\smallbreak
		On en déduit $\boxed{\Pc_L=-B_0ia\cos(\alpha)\xp}$~: \fbox{\textbf{réponse
				C)}}
		\nitem{5} On étudie la barre en translation selon $\ex$, dans le référentiel
		terrestre supposé galiléen.
		\smallbreak
		La barre est soumise au poids, à la résultante des forces de Laplace et
		aux 2 réactions au niveau des rails (forces orthogonales à $\ex$ en
		l'absence de frottement de glissement).
		\figsvg{figcorr.pdf_tex}
		On applique le PFD projeté selon $\ex$ pour obtenir la \fbox{\textbf{réponse
				A)}}~:
		\[
			\boxed{m\xpp=mg\sin(\alpha)-B_0ia\cos(\alpha)}
		\]
		\nitem{6} On utilise l'expression de $i(t)=I_0\cos(\wt)$.
		\smallbreak
		Par intégration avec la CI $\xp(t=0)=0$~:
		\[
			\xp(t)=g\sin(\alpha)t-\cfrac{B_0I_0a\cos(\alpha)}{m\w}\sin(\wt)
		\]
		On intègre encore une fois avec la CI $x(t=0)=x_0$ pour obtenir la
		\fbox{\textbf{réponse D)}}~:
		\[
			\boxed{x(t)=\cfrac{g\sin(\alpha)}{2}t^2+
			\cfrac{B_0I_0a\cos(\alpha)}{m\w^2}(\cos(\wt)-1)+x_0}
		\]
	\end{enumerate}
}

\end{document}
